\Article{24}
The triangles among five points of the Hanover triangulation—Falkenberg, Breithorn, Hauselberg, Wulfsode, and Wilsede—will furnish a second example. The directions were observed\footnote{The reference directions to which the individual directions are referred are here regarded as arbitrary, although in fact they coincide with the meridian lines of the stations. The observations will be published in full later; for the present, the figure is found in \textit{Astronomische Nachrichten}, Vol. I, p. 441.}:
\begin{center}
\begin{tabular}{r l l r r r}
\multicolumn{6}{c}{At station \textit{Falkenberg}}\\[5pt]
0.&Wilsede&. . . . .&\(187^0\)&\(47'\)&\(30''311\)\\
1.&Wulfsode&. . . .&225&9&39,676\\
2.&Hauselberg&. . .&266&13&56,239\\
3.&Breithorn&. . . .&274&14&43,634
\end{tabular}
\end{center}
\begin{center}
\begin{tabular}{r l l r r r}
\multicolumn{6}{c}{At station \textit{Breithorn}}\\[5pt]
4.&Falkenberg&. . .&94&33&40,755\\
5.&Hauselberg&. . .&122&51&23,054\\
6.&Wilsede&. . . . .&150&18&35,100
\end{tabular}
\end{center}
\begin{center}
\begin{tabular}{r l l r r r}
\multicolumn{6}{c}{At station \textit{Hauselberg}}\\[5pt]
7.&Falkenberg&. . .&86&29&6,872\\
8.&Wilsede&. . . . .&154&37&9,624\\
9.&Wulfsode&. . . .&189&2&56,376\\
10.&Breithorn&. . . .&302&47&37,732
\end{tabular}
\end{center}
\begin{center}
\begin{tabular}{r l l r r r}
\multicolumn{6}{c}{At station \textit{Wulfsode}}\\[5pt]
11.&Hauselberg&. . .&9&5&36,593\\
12.&Falkenberg&. . .&45&27&33,556\\
13.&Wilsede&. . . . .&118&44&13,159
\end{tabular}
\end{center}
\SourceLeaf{093}
\begin{center}
\begin{tabular}{r l l r r r}
\multicolumn{6}{c}{At station \textit{Wilsede}}\\[5pt]
14.&Falkenberg&. . .&\(7^0\)&\(51'\)&\(1''027\)\\
15.&Wulfsode&. . . .&298&29&49,519\\
16.&Breithorn&. . . .&330&3&7,392\\
17.&Hauselberg&. . .&334&25&26,746
\end{tabular}
\end{center}

These observations allow seven triangles to be formed.
\begin{center}
\begin{tabular}{l l r r r}
\multicolumn{5}{c}{Triangle I.}\\[5pt]
Falkenberg&. . . . . .&\(8^0\)&\(0'\)&\(47''395\)\\
Breithorn&. . . . . . .&28&17&42,299\\
Hauselberg&. . . . . .&143&41&29,140
\end{tabular}
\end{center}
\begin{center}
\begin{tabular}{l l r r r}
\multicolumn{5}{c}{Triangle II.}\\[5pt]
Falkenberg&. . . . . .&86&27&13,323\\
Breithorn&. . . . . . .&55&44&54,345\\
Wilsede&. . . . . . . .&37&47&53.635
\end{tabular}
\end{center}
\begin{center}
\begin{tabular}{l l r r r}
\multicolumn{5}{c}{Triangle III.}\\[5pt]
Falkenberg&. . . . . .&41&4&16,563\\
Hauselberg&. . . . . .&102&33&49,504\\
Wulfsode&. . . . . . .&36&21&56,963
\end{tabular}
\end{center}
\begin{center}
\begin{tabular}{l l r r r}
\multicolumn{5}{c}{Triangle IV.}\\[5pt]
Falkenberg&. . . . . .&78&26&25,928\\
Hauselberg&. . . . . .&68&8&2,752\\
Wilsede&. . . . . . . .&35&25&34,281
\end{tabular}
\end{center}
\begin{center}
\begin{tabular}{l l r r r}
\multicolumn{5}{c}{Triangle V.}\\[5pt]
Falkenberg&. . . . . .&37&22&9,365\\
Wulfsode&. . . . . . .&73&16&39,603\\
Wilsede&. . . . . . . .&69&21&11,508
\end{tabular}
\end{center}
\begin{center}
\begin{tabular}{l l r r r}
\multicolumn{5}{c}{Triangle VI.}\\[5pt]
Breithorn&. . . . . . .&27&27&12,046\\
Hauselberg&. . . . \hspace{.5em} .&148&10&28,108\\
Wilsede&. . . . . . . .&4&22&19,354
\end{tabular}
\end{center}
\SourceLeaf{094}
\begin{center}
\begin{tabular}{l l r r r}
\multicolumn{5}{c}{Triangle VII.}\\[5pt]
Hauselberg&. . . . . .&\(34^0\)&\(25'\)&\(46''752\)\\
Wulfsode&. . . . . . .&109&38&36,566\\
Wilsede&. . . . . . . .&35&55&37,227
\end{tabular}
\end{center}

There will therefore be seven conditional equations of the second kind (equations of the first kind are evidently absent). To derive them, we must first compute the spheroidal excesses of the seven triangles. For this purpose, the absolute magnitude of at least one side must be known: the side between Wilsede and Wulfsode is 22377,94 metres. The spheroidal excesses of triangles I, II, III, IV, V, VI, and VII are therefore, in order: I\(.\,.\,.\;0''202\); II\(.\,.\,.\;2''442\); III\(.\,.\,.\;1''257\); IV\(.\,.\,.\;1''919\); V\(.\,.\,.\;1''957\); VI\(.\,.\,.\;0''321\); VII\(.\,.\,.\;1''295\).

Now, if the directions are denoted by \(v^{(0)},v^{(1)},v^{(2)},v^{(3)}\) etc., in the order given above and distinguished by the indices, the angles of triangle I are
\[
v^{(3)}-v^{(2)},\quad v^{(5)}-v^{(4)},\quad 360^0+v^{(7)}-v^{(10)}
\]
and hence the first conditional equation is
\[
-v^{(2)}+v^{(3)}-v^{(4)}+v^{(5)}+v^{(7)}-v^{(10)}+179^0\,59'\,59''798=0
\]
The other six triangles likewise furnish further equations. A little attention, however, shows that these seven equations are not independent: the second is identical with the sum of the first, fourth, and sixth, and the sum of the third and fifth is identical with the sum of the fourth and seventh. We shall therefore omit the second and fifth. In place of the remaining conditional equations in finite form, we write the corresponding equations from system (13), using (0), (1), (2), etc. for the symbols \(\varepsilon,\varepsilon'\), etc.:
\[
\begin{aligned}
-1''368&=-(2)+(3)-(4)+(5)+(7)-(10)\\
+1{,}773&=-(1)+(2)-(7)+(9)-(11)+(12)\\
+1{,}042&=-(0)+(2)-(7)+(8)+(14)-(17)\\
-0{,}813&=-(5)+(6)-(8)+(10)-(16)+(17)\\
-0{,}750&=-(8)+(9)-(11)+(13)-(15)+(17)
\end{aligned}
\]

\textit{Eight} conditional equations of the third kind could be obtained from the system of triangles, since triples could be formed both from the four triangles I, II, IV, and VI and from the following
\SourceLeaf{095}
III, IV, V, and VII for this purpose. A little attention nevertheless shows that \textit{two} suffice, one from each group, since the others must already be contained in these and the preceding conditional equations. The sixth conditional equation will therefore be
\[
\begin{aligned}
&\log\sin(v^{(3)}-v^{(2)}-0''067)-\log\sin(v^{(5)}-v^{(4)}-0''067)\\
&+\log\sin(v^{(14)}-v^{(17)}-0''640)-\log\sin(v^{(2)}-v^{(0)}-0''640)\\
&+\log\sin(v^{(6)}-v^{(5)}-0''107)-\log\sin(v^{(17)}-v^{(16)}-0''107)=0
\end{aligned}
\]
and the seventh
\[
\begin{aligned}
&\log\sin(v^{(2)}-v^{(1)}-0''419)-\log\sin(v^{(12)}-v^{(11)}-0''419)\\
&+\log\sin(v^{(14)}-v^{(17)}-0''640)-\log\sin(v^{(2)}-v^{(0)}-0''640)\\
&+\log\sin(v^{(13)}-v^{(11)}-0''432)-\log\sin(v^{(17)}-v^{(15)}-0''432)=0
\end{aligned}
\]
with the corresponding equations from system (13) being
\[
\begin{aligned}
+25={}&+4{,}31(0)-153{,}88(2)+149{,}57(3)+39{,}11(4)-79{,}64(5)\\
&+40{,}53(6)+31{,}90(14)+275{,}39(16)-307{,}29(17)\\
-3={}&+4{,}31(0)-24{,}16(1)+19{,}85(2)+36{,}11(11)-28{,}59(12)\\
&-7{,}52(13)+31{,}90(14)+29{,}06(15)-60{,}96(17)
\end{aligned}
\]

If we now assign the same degree of certainty to each direction, setting \(p^{(0)}=p^{(1)}=p^{(2)}\ \text{etc.}=1\), and denote by \(A,B,C,D,E,F,G\) the correlates of the seven conditional equations in the order followed here, their determination is to be obtained from the following equations:
\[
\begin{aligned}
-1{,}368={}&+6A-2B-2C-2D+184{,}72F-19{,}85G\\
+1{,}773={}&-2A+6B+2C+2E-153{,}88F-20{,}69G\\
+1{,}042={}&-2A+2B+6C-2D-2E+181{,}00F+108{,}40G\\
-0{,}813={}&-2A-2C+6D+2E-462{,}51F-60{,}96G\\
-0{,}750={}&+2B-2C+2D+6E-307{,}29F-133{,}65G\\
+25={}&+184{,}72A-153{,}88B+181{,}00C-462{,}51D-307{,}29E\\
&+224868F+16694{,}1G\\
-3={}&-19{,}85A-20{,}69B+108{,}40C-60{,}96D-133{,}65E\\
&+16694{,}1F+8752{,}39G
\end{aligned}
\]

By elimination we deduce
\SourceLeaf{096}
\[
\begin{aligned}
A&=-0{,}225\\
B&=+0{,}344\\
C&=-0{,}088\\
D&=-0{,}171\\
E&=-0{,}323\\
F&=+0{,}000215915\\
G&=-0{,}00547462
\end{aligned}
\]

The most probable errors are now given by the formulas:
\[
\begin{aligned}
(0)&=-C+4{,}31F+4{,}31G\\
(1)&=-B-24{,}16G\\
(2)&=-A+B+C-153{,}88F+19{,}85G
\end{aligned}
\]
etc., from which the numerical values follow
\[
\begin{array}{r@{\;}l@{\qquad}|@{\qquad}r@{\;}l}
(0)=&+0''065&(9)=&+0''021\\
(1)=&-0{,}212&(10)=&+0{,}054\\
(2)=&+0{,}339&(11)=&-0{,}219\\
(3)=&-0{,}193&(12)=&+0{,}501\\
(4)=&+0{,}233&(13)=&-0{,}282\\
(5)=&-0{,}071&(14)=&-0{,}256\\
(6)=&-0{,}162&(15)=&+0{,}164\\
(7)=&-0{,}481&(16)=&+0{,}230\\
(8)=&+0{,}406&(17)=&-0{,}139
\end{array}
\]

The sum of the squares of these errors is found to be \(=1{,}2288\); hence the mean error of one direction, insofar as it can be inferred from the 18 observed directions, is
\[
=\BareRoot{\frac{1{,}2288}{7}}=0''4190
\]

